\section{Class PSPACE --- TQBFHard Semantics}
%%% generated by the blueprint pipeline from TCSlib.Complexity.ClassPSPACE.TQBFHardSemantics
\subsection*{Class PSPACE --- TQBFHard Semantics: Introduction}

Declarations of the Lean module \texttt{TCSlib.Complexity.ClassPSPACE.TQBFHardSemantics}.

\subsection*{Class PSPACE --- TQBFHard Semantics: Results}

\begin{definition}[Writing a Boolean block]
\lean{Complexity.TQBFHard.writeBlock}
\label{Complexity.TQBFHard.writeBlock}
\leanok
Given a Boolean sequence $\sigma$, a starting index $i$, and a finite Boolean list
$\beta$, define a new Boolean sequence by replacing the entries from index $i$ through
index $i+|\beta|-1$ with the entries of $\beta$ in order, and leaving all other entries
equal to $\sigma$.
\end{definition}

\begin{definition}[Reading a Boolean block]
\lean{Complexity.TQBFHard.readBlock}
\label{Complexity.TQBFHard.readBlock}
\leanok
Given a Boolean sequence $\tau$, a starting index $i$, and a length $w$, the block read
from $i$ is the length-$w$ list whose entry at position $j$, for $0\le j<w$, is
$\tau(i+j)$.
\end{definition}

\begin{theorem}[Length of a read block]
\lean{Complexity.TQBFHard.length_readBlock}
\label{Complexity.TQBFHard.length_readBlock}
\leanok
\uses{Complexity.TQBFHard.readBlock}
\difficulty{1}
Let $\tau:\mathbb{N}\to\{\mathrm{false},\mathrm{true}\}$ be a bit sequence, and let
$i,w\in\mathbb{N}$. The block consisting of the $w$ bits
$\tau(i),\tau(i+1),\ldots,\tau(i+w-1)$ has length $w$.
\end{theorem}

\begin{theorem}[Writing before the block]
\lean{Complexity.TQBFHard.writeBlock_of_lt}
\label{Complexity.TQBFHard.writeBlock_of_lt}
\leanok
\uses{Complexity.TQBFHard.writeBlock}
\difficulty{3}
Let $\sigma:\mathbb{N}\to\{\mathrm{false},\mathrm{true}\}$ be a bit sequence, let
$i\in\mathbb{N}$, and let $\beta$ be a finite bit string. At every position $v<i$,
writing $\beta$ starting at position $i$ leaves the bit at $v$ equal to $\sigma(v)$.
\end{theorem}

\begin{theorem}[Reading a written block]
\lean{Complexity.TQBFHard.writeBlock_of_mem}
\label{Complexity.TQBFHard.writeBlock_of_mem}
\leanok
\uses{Complexity.TQBFHard.writeBlock}
\difficulty{3}
Let $\sigma:\mathbb{N}\to\{\mathrm{false},\mathrm{true}\}$ be a bit sequence, let
$i\in\mathbb{N}$, and let $\beta$ be a finite bit string. If $j<|\beta|$, then after
writing $\beta$ starting at position $i$, the bit at position $i+j$ is the $j$th bit of
$\beta$ (with default value false if that position is out of range).
\end{theorem}

\begin{theorem}[Prepending a bit to a written block]
\lean{Complexity.TQBFHard.writeBlock_update_cons}
\label{Complexity.TQBFHard.writeBlock_update_cons}
\leanok
\uses{Complexity.TQBFHard.writeBlock}
\difficulty{5}
Let $\sigma:\mathbb{N}\to\{\mathrm{false},\mathrm{true}\}$ be a bit sequence, let
$i\in\mathbb{N}$, let $b$ be a bit, and let $\beta$ be a finite bit string. Writing
$\beta$ starting at position $i+1$ in the sequence obtained from $\sigma$ by setting its
bit at position $i$ to $b$ gives the same sequence as writing $b$ followed by $\beta$
starting at position $i$.
\end{theorem}

\begin{theorem}[Writing an empty block]
\lean{Complexity.TQBFHard.writeBlock_nil}
\label{Complexity.TQBFHard.writeBlock_nil}
\leanok
\uses{Complexity.TQBFHard.writeBlock}
\difficulty{3}
Let $\sigma:\mathbb{N}\to\{\mathrm{false},\mathrm{true}\}$ be a bit sequence and let
$i\in\mathbb{N}$. Writing an empty block into $\sigma$ starting at position $i$ leaves
the sequence unchanged.
\end{theorem}

\begin{theorem}[A read block equals a prescribed word]
\lean{Complexity.TQBFHard.readBlock_eq}
\label{Complexity.TQBFHard.readBlock_eq}
\leanok
\uses{Complexity.TQBFHard.readBlock}
\difficulty{3}
Let $\tau:\mathbb N\to\{0,1\}$, $i\in\mathbb N$, and let $\beta$ be a finite bit string.
If $\tau(i+j)$ equals the $j$th bit of $\beta$ for every $0\le j<|\beta|$, then the
block of $|\beta|$ bits read from position $i$ in $\tau$ is $\beta$.
\end{theorem}

\begin{theorem}[A read block equals a word of matching length]
\lean{Complexity.TQBFHard.readBlock_eq'}
\label{Complexity.TQBFHard.readBlock_eq'}
\leanok
\uses{Complexity.TQBFHard.readBlock, Complexity.TQBFHard.readBlock_eq}
\difficulty{1}
Let $\tau:\mathbb N\to\{0,1\}$, let $i,w\in\mathbb N$, and let $\beta$ be a finite bit
string of length $w$. If $\tau(i+j)$ equals the $j$th bit of $\beta$ for every $0\le
j<w$, then the block of $w$ bits read from position $i$ in $\tau$ is $\beta$.
\end{theorem}

\begin{theorem}[Agreement determines a read block]
\lean{Complexity.TQBFHard.readBlock_congr}
\label{Complexity.TQBFHard.readBlock_congr}
\leanok
\uses{Complexity.TQBFHard.readBlock}
\difficulty{3}
Let $\tau,\tau':\mathbb N\to\{0,1\}$ and let $i,w\in\mathbb N$. If
$\tau(i+j)=\tau'(i+j)$ for every integer $j$ with $0\le j<w$, then the length-$w$ blocks
read from position $i$ in these two sequences are equal.
\end{theorem}

\begin{theorem}[Existential Prefix Block Compression]
\lean{Complexity.TQBFHard.prefixTruth_replicate_ex}
\label{Complexity.TQBFHard.prefixTruth_replicate_ex}
\leanok
\uses{Complexity.QBF.PrefixTruth, Complexity.QBF.Quant, Complexity.TQBFHard.writeBlock, Complexity.TQBFHard.writeBlock_update_cons}
\difficulty{6}
Let $P$ be a property of infinite Boolean assignments, let $w,i\in\mathbb{N}$, let
$\sigma:\mathbb{N}\to\{\mathrm{false},\mathrm{true}\}$ be an assignment, and let
$\mathrm{rest}$ be a finite sequence of existential or universal quantifiers. The prefix
truth obtained by prepending $w$ existential quantifiers to $\mathrm{rest}$ is
equivalent to the existence of a Boolean word $\beta$ of length $w$ such that the
remaining prefix is true after replacing the assignment on positions $i,\ldots,i+w-1$ by
$\beta$ and starting the remaining prefix at position $i+w$.
\end{theorem}

\begin{theorem}[Universal Prefix Block Compression]
\lean{Complexity.TQBFHard.prefixTruth_replicate_all}
\label{Complexity.TQBFHard.prefixTruth_replicate_all}
\leanok
\uses{Complexity.QBF.PrefixTruth, Complexity.QBF.Quant, Complexity.TQBFHard.writeBlock, Complexity.TQBFHard.writeBlock_update_cons}
\difficulty{6}
Let $P$ be a property of infinite Boolean assignments, let $w,i\in\mathbb{N}$, let
$\sigma:\mathbb{N}\to\{\mathrm{false},\mathrm{true}\}$ be an assignment, and let
$\mathrm{rest}$ be a finite sequence of existential or universal quantifiers. The prefix
truth obtained by prepending $w$ universal quantifiers to $\mathrm{rest}$ is equivalent
to requiring, for every Boolean word $\beta$ of length $w$, that $P$ with the remaining
prefix be true after replacing the assignment on positions $i,\ldots,i+w-1$ by $\beta$
and starting the remaining prefix at position $i+w$.
\end{theorem}

\begin{theorem}[Implication through a quantified prefix]
\lean{Complexity.TQBFHard.prefixTruth_imp_const}
\label{Complexity.TQBFHard.prefixTruth_imp_const}
\leanok
\uses{Complexity.QBF.PrefixTruth, Complexity.QBF.Quant}
\difficulty{5}
Let $K$ be a proposition, let $P$ be a predicate on Boolean assignments
$\tau:\mathbb{N}\to\{\mathrm{false},\mathrm{true}\}$, let $qs$ be a finite list of
existential and universal quantifiers, and let $\sigma$ be a Boolean assignment with
starting index $i$. The quantified formula with matrix $K\to P(\tau)$ is true under
$qs,\sigma,i$ if and only if $K$ implies that the quantified formula with matrix
$P(\tau)$ is true under the same prefix, assignment, and index.
\end{theorem}

\begin{theorem}[Conjunction through a quantified prefix]
\lean{Complexity.TQBFHard.prefixTruth_and_const}
\label{Complexity.TQBFHard.prefixTruth_and_const}
\leanok
\uses{Complexity.QBF.PrefixTruth, Complexity.QBF.Quant}
\difficulty{5}
Let $K$ be a proposition, let $P$ be a predicate on Boolean assignments
$\tau:\mathbb{N}\to\{\mathrm{false},\mathrm{true}\}$, let $qs$ be a finite list of
existential and universal quantifiers, and let $\sigma$ be a Boolean assignment with
starting index $i$. The quantified formula with matrix $K\land P(\tau)$ is true under
$qs,\sigma,i$ if and only if $K$ holds and the quantified formula with matrix $P(\tau)$
is true under the same prefix, assignment, and index.
\end{theorem}

\begin{definition}[Quantifier list by level]
\lean{Complexity.TQBFHard.levelQuants}
\label{Complexity.TQBFHard.levelQuants}
\leanok
\uses{Complexity.QBF.Quant}
For each width $w\in\mathbb{N}$, define a sequence $Q_w$ of lists of quantifiers by
$Q_w(0)=[]$ and, for every $r\in\mathbb{N}$, let $Q_w(r+1)$ be the concatenation of $w$
existential quantifiers, $2w$ universal quantifiers, and $Q_w(r)$, in that order.
\end{definition}

\begin{theorem}[Length of the level quantifier list]
\lean{Complexity.TQBFHard.length_levelQuants}
\label{Complexity.TQBFHard.length_levelQuants}
\leanok
\uses{Complexity.TQBFHard.levelQuants}
\difficulty{3}
For natural numbers $w$ and $r$, the list of quantifiers at width $w$ and depth $r$ has
length $3wr$.
\end{theorem}

\begin{definition}[Boolean list selection relation]
\lean{Complexity.TQBFHard.Sel}
\label{Complexity.TQBFHard.Sel}
\leanok
For Boolean lists $A$, $B$, $z$, $u$, and $v$, the selection relation holds precisely
when either $u=A$ and $v=z$, or $u=z$ and $v=B$.
\end{definition}

\begin{definition}[Recursive level matrix condition]
\lean{Complexity.TQBFHard.levelMatrix}
\label{Complexity.TQBFHard.levelMatrix}
\leanok
\uses{Complexity.TQBFHard.Sel, Complexity.TQBFHard.readBlock}
Let $w,r,p\in\mathbb{N}$, let $E$ be a relation on Boolean lists, let $A,B$ be Boolean
lists, and let $\tau:\mathbb{N}\to\{\mathrm{false},\mathrm{true}\}$. For
$i\in\mathbb{N}$, write $R_i$ for the length-$w$ block
$(\tau(i),\tau(i+1),\ldots,\tau(i+w-1))$. At level $0$, the condition is $E(A,B)$. At
level $r+1$, the condition is that if either $R_p=A$ and $R_{p+w}=R_{p+2w}$, or
$R_p=R_{p+2w}$ and $R_{p+w}=B$, then the level-$r$ condition holds at index $p+3w$, with
endpoints $R_{p+w}$ and $R_{p+2w}$ and sequence $\tau$.
\end{definition}

\begin{definition}[Recursive midpoint relation]
\lean{Complexity.TQBFHard.Mid}
\label{Complexity.TQBFHard.Mid}
\leanok
Fix a width $w$ and a relation $E$ on pairs of Boolean lists. Define a relation
recursively on the level and endpoint lists: at level zero it is $E(A,B)$; at level
$r+1$, it holds when there is a Boolean list $z$ of length $w$ such that the level-$r$
relation holds from $A$ to $z$ and from $z$ to $B$.
\end{definition}

\begin{theorem}[Quantified level matrix and midpoint relation]
\lean{Complexity.TQBFHard.prefixTruth_levelMatrix}
\label{Complexity.TQBFHard.prefixTruth_levelMatrix}
\leanok
\uses{Complexity.QBF.PrefixTruth, Complexity.QBF.prefixTruth_congr_agree, Complexity.TQBFHard.Mid, Complexity.TQBFHard.Sel, Complexity.TQBFHard.levelMatrix, Complexity.TQBFHard.levelQuants, Complexity.TQBFHard.prefixTruth_imp_const, Complexity.TQBFHard.prefixTruth_replicate_all, Complexity.TQBFHard.prefixTruth_replicate_ex, Complexity.TQBFHard.readBlock, Complexity.TQBFHard.readBlock_eq', Complexity.TQBFHard.writeBlock, Complexity.TQBFHard.writeBlock_of_lt, Complexity.TQBFHard.writeBlock_of_mem}
\difficulty{6}
Let $w,r$ be natural numbers and let $E$ be a relation on pairs of finite Boolean
strings. Define $\operatorname{Mid}_0(A,B)$ to mean $E(A,B)$; define
$\operatorname{Mid}_{r+1}(A,B)$ to mean that some string $z$ of length $w$ satisfies
both $\operatorname{Mid}_r(A,z)$ and $\operatorname{Mid}_r(z,B)$. For every position
$p$, Boolean assignment $\sigma$, and strings $A,B$ of length $w$, the quantified
formula with prefix consisting of $r$ repetitions of $w$ existential, $w$ universal, and
$w$ universal quantifiers, and with matrix given by the recursively defined level matrix
starting at position $p$ with endpoints $A,B$, is true under assignment $\sigma$ and
starting index $p$ if and only if $\operatorname{Mid}_r(A,B)$ holds.
\end{theorem}

\begin{theorem}[Top-Level Prefix Characterization]
\lean{Complexity.TQBFHard.prefixTruth_top}
\label{Complexity.TQBFHard.prefixTruth_top}
\leanok
\uses{Complexity.QBF.PrefixTruth, Complexity.QBF.prefixTruth_congr_agree, Complexity.TQBFHard.Mid, Complexity.TQBFHard.levelMatrix, Complexity.TQBFHard.levelQuants, Complexity.TQBFHard.prefixTruth_and_const, Complexity.TQBFHard.prefixTruth_levelMatrix, Complexity.TQBFHard.prefixTruth_replicate_ex, Complexity.TQBFHard.readBlock, Complexity.TQBFHard.readBlock_eq', Complexity.TQBFHard.writeBlock, Complexity.TQBFHard.writeBlock_of_lt, Complexity.TQBFHard.writeBlock_of_mem}
\difficulty{5}
Let $w,R\in\mathbb{N}$, let $E$ be a relation on finite Boolean words, let $I$ and $J$
be predicates on such words, and let $\sigma$ be an infinite Boolean assignment. The
quantified formula with three consecutive blocks of $w$ existential quantifiers followed
by $R$ levels of the prefix specified by the recursive pattern “existential block,
universal block, universal block” is true from position $0$, with matrix requiring $I$
on the block read at positions $w$ through $2w-1$, $J$ on the block read at positions
$2w$ through $3w-1$, and the recursively defined level matrix from position $3w$, if and
only if there are Boolean words $A,B$ of length $w$ for which $I(A)$ and $J(B)$ hold and
$\mathrm{Mid}_w(E,R,A,B)$ holds. Here $\mathrm{Mid}_w(E,0,A,B)$ means $E(A,B)$, and
$\mathrm{Mid}_w(E,r+1,A,B)$ means that some word $z$ of length $w$ satisfies both
$\mathrm{Mid}_w(E,r,A,z)$ and $\mathrm{Mid}_w(E,r,z,B)$.
\end{theorem}

\begin{definition}[Recursive relation for a fixed number of steps]
\lean{Complexity.TQBFHard.Steps}
\label{Complexity.TQBFHard.Steps}
\leanok
Fix a width $w$ and a relation $E$ on pairs of Boolean lists. At zero steps, the
relation between $A$ and $B$ means $A=B$. At $k+1$ steps, it means there is a Boolean
list $z$ of length $w$ such that the relation holds from $A$ to $z$ in $k$ steps and
$E(z,B)$ holds.
\end{definition}

\begin{theorem}[Addition of step counts]
\lean{Complexity.TQBFHard.steps_add}
\label{Complexity.TQBFHard.steps_add}
\leanok
\uses{Complexity.TQBFHard.Steps}
\difficulty{4}
For a width $w$, a relation $E$ on Boolean strings, and a natural number $r$, define an
$r$-step chain from $A$ to $B$ by requiring $A=B$ when $r=0$, and otherwise requiring a
length-$w$ string $z$ with an $(r-1)$-step chain from $A$ to $z$ and $E(z,B)$. For all
natural numbers $m,k$ and strings $A,B$ of length $w$, an $(k+m)$-step chain from $A$ to
$B$ exists if and only if there is a length-$w$ string $z$ such that there is a $k$-step
chain from $A$ to $z$ and an $m$-step chain from $z$ to $B$.
\end{theorem}

\begin{theorem}[Midpoint relation and bounded step chains]
\lean{Complexity.TQBFHard.mid_iff_steps}
\label{Complexity.TQBFHard.mid_iff_steps}
\leanok
\uses{Complexity.TQBFHard.Mid, Complexity.TQBFHard.Steps, Complexity.TQBFHard.steps_add}
\difficulty{4}
Let $w,r\in\mathbb{N}$ and let $E$ be a relation on finite bit strings. For all bit
strings $A,B$ of length $w$, the midpoint relation of depth $r$ holds between $A$ and
$B$ if and only if there is a chain of $2^r$ transitions from $A$ to $B$ in which every
intermediate string has length $w$ and each transition satisfies $E$.
\end{theorem}

\begin{definition}[First block of the encoding]
\lean{Complexity.TQBFHard.zRec}
\label{Complexity.TQBFHard.zRec}
\leanok
\uses{Complexity.TQBFHard.readBlock}
For a Boolean sequence $\tau : \mathbb{N} \to \{\mathrm{false},\mathrm{true}\}$ and
natural numbers $w$ and $j$, let $z(\tau,w,j)$ be the length-$w$ block of values of
$\tau$ beginning at position $3wj$.
\end{definition}

\begin{definition}[Second block of the encoding]
\lean{Complexity.TQBFHard.uRec}
\label{Complexity.TQBFHard.uRec}
\leanok
\uses{Complexity.TQBFHard.readBlock}
For a Boolean sequence $\tau : \mathbb{N} \to \{\mathrm{false},\mathrm{true}\}$ and
natural numbers $w$ and $j$, let $u(\tau,w,j)$ be the length-$w$ block of values of
$\tau$ beginning at position $3wj+w$.
\end{definition}

\begin{definition}[Third block of the encoding]
\lean{Complexity.TQBFHard.vRec}
\label{Complexity.TQBFHard.vRec}
\leanok
\uses{Complexity.TQBFHard.readBlock}
For a Boolean sequence $\tau : \mathbb{N} \to \{\mathrm{false},\mathrm{true}\}$ and
natural numbers $w$ and $j$, let $v(\tau,w,j)$ be the length-$w$ block of values of
$\tau$ beginning at position $3wj+2w$.
\end{definition}

\begin{theorem}[Flattened characterization of the level matrix]
\lean{Complexity.TQBFHard.levelMatrix_iff_flat}
\label{Complexity.TQBFHard.levelMatrix_iff_flat}
\leanok
\uses{Complexity.TQBFHard.Sel, Complexity.TQBFHard.levelMatrix, Complexity.TQBFHard.uRec, Complexity.TQBFHard.vRec, Complexity.TQBFHard.zRec}
\difficulty{5}
Let $w,r\in\mathbb{N}$, let $E$ be a relation on pairs of finite bit strings, and let
$\tau:\mathbb{N}\to\{\mathrm{false},\mathrm{true}\}$ be a bit sequence. For every
$j\in\mathbb{N}$, the recursively defined level-matrix condition at position $3w(j+1)$,
with endpoint strings $u_j$ and $v_j$ given by the length-$w$ blocks of $\tau$ starting
at $3wj+w$ and $3wj+2w$, respectively, holds if and only if either at some index $i<r$
the selector condition fails for the five blocks
$u_{j+i},v_{j+i},z_{j+i+1},u_{j+i+1},v_{j+i+1}$, or $E(u_{j+r},v_{j+r})$ holds. Here
$z_k$ is the length-$w$ block starting at $3wk$, and the selector condition means that
either $(u_{j+i},v_{j+i})=(u_{j+i+1},z_{j+i+1})$ or
$(u_{j+i+1},v_{j+i+1})=(z_{j+i+1},v_{j+i})$.
\end{theorem}
